% Generated draft: requires mathematical and source review.
\documentclass[11pt]{article}
\usepackage[margin=25mm]{geometry}
\usepackage{fontspec}
\setmainfont{Latin Modern Roman}
\usepackage{amsmath,amssymb,mathtools,unicode-math}
\setmathfont{Latin Modern Math}
\usepackage{enumitem,needspace,xurl,hyperref}
\setlength{\parindent}{0pt}
\setlength{\parskip}{5pt}
\setlength{\emergencystretch}{3em}
\newcommand{\N}{\mathbb{N}}
\newcommand{\Z}{\mathbb{Z}}
\newcommand{\Q}{\mathbb{Q}}
\newcommand{\R}{\mathbb{R}}
\newcommand{\C}{\mathbb{C}}
\newcommand{\F}{\mathbb{F}}
\newcommand{\A}{\mathbb{A}}
\newcommand{\PP}{\mathbb{P}}
\newcommand{\E}{\mathbb{E}}
\newcommand{\eps}{\varepsilon}
\newcommand{\dd}{\,\mathrm{d}}
\begin{document}
\section*{UnsolvedMath ID: 10900038}
\textbf{Source problem code:} catalogue entry
\subsection{Definitions and mathematical statement}
Let $F_r$ be the free group on $r$ generators, where $r$ is an integer with $r\ge2$. For a group $G$, define
\[
 \mathcal C(G)=\{[Q]: Q\text{ is a finite group and }G\twoheadrightarrow Q\},
\]
the set of isomorphism classes of its finite quotient groups. This is a set of quotient types, without counts of epimorphisms or a choice of quotient maps, as stipulated in source item 5.12.

\textbf{Question (Reid, item 5.13).} Does $\mathcal C(F_r)$ determine $F_r$ up to isomorphism? The comparison class of groups is not stated in the extracted question. The research summary identifies the intended general conjecture with the following finitely generated residually finite formulation:
\[
 \forall r\ge2\ \forall G\quad
 \bigl(G\text{ finitely generated and residually finite},\ 
       \mathcal C(G)=\mathcal C(F_r)\bigr)
 \ \Longrightarrow\ G\cong F_r.
\]
Here residually finite means that every nonidentity element survives in some finite quotient. If
\[
 \widehat G=\varprojlim_{K\triangleleft G,\,[G:K]<\infty}G/K
\]
is the profinite completion, the corresponding usual formulation for this finitely generated comparison class asks whether $\widehat G\cong\widehat{F_r}$ as topological groups forces $G\cong F_r$. The summary calls this Remeslennikov's conjecture.

The finitely generated and residually finite hypotheses come from the supplied summary, not from the one-line source excerpt; they must be confirmed as the intended comparison class. Finite quotients alone do not retain elements in the kernel of the natural map to the profinite completion. The question is not merely whether two finite-rank free groups have the same rank.
\subsection{Short English statement}
If a finitely generated residually finite group has exactly the same finite quotient types as a free group of rank at least two, must it be that free group? This is the comparison class specified by the supplied summary; the original excerpt leaves it implicit.
\subsection{Sources}
\begin{itemize}
\item Local UnsolvedMath registry, permanent ID 10900038, source code catalogue entry; attributed to Reid, no proposal year supplied. The definition of $\mathcal C$ is inherited from item 5.12, also preserved in local permanent ID 10900037.
\item Delp, Hoffoss and Manning, \emph{Problems In Groups, Geometry, and Three-Manifolds} (2015), subsection 5.13 (Reid), PDF page 4; \url{https://arxiv.org/pdf/1512.04620#page=4} and \url{https://arxiv.org/abs/1512.04620}. The registry records source-TeX extraction checked against rendered PDF text and access on 2026-07-29.
\item The research summary names Remeslennikov's conjecture and cites \emph{On the profinite rigidity of free and surface groups}, \url{https://arxiv.org/abs/2211.12390}, for a restricted class described as $\mathcal X_{\mathrm{ab}}$ with a finite abelian hierarchy starting from residually free groups.
\item The summary also cites \emph{Profinite detection of free products and free factors}, \url{https://arxiv.org/abs/2603.16674} (2026), as additional progress. The supplied titles, identifiers, and claimed theorem scopes are not independently checked here.
\end{itemize}
\end{document}
